%! TEX root = ../am2-20-21.tex

\chapter{Serii Fourier}

Pentru funcțiile care au proprietăți de periodicitate (sau pot fi transformate
în unele cu astfel de proprietăți), ne interesează o dezvoltare specială a lor,
anume în serii de sinusuri și cosinusuri, numită \emph{serie Fourier}.

Mai exact, în general, pentru o funcție cu proprietățile potrivite, vom avea
o dezvoltare de forma:
\[
    f(x) = \sum_{n \geq 0} \left[ a_n \cos(nx) + b_n \sin(nx) \right]
\]

Sumele parțiale ale acestei serii se numesc polinoame trigonometrice,
notate în general $ P_n(x) $ și definite prin:
\[
    P_n(x) = \sum_{k = 0}^n a_k \cos(kx) + b_k \sin(kx).
\]

Problema principală de interes este determinarea coeficienților
care apar în seria Fourier. Vom ține cont de identitățile de mai jos,
datorate parității funcțiilor trigonometrice:
\begin{align*}
    \ds\int_0^{2\pi} \cos kx dx &= \begin{cases} 0, & k \neq 0 \\ 2 \pi, & k = 0 \end{cases} \\
    \ds\int_0^{2\pi} \sin kx dx &= \begin{cases} 0, & k \neq 0 \\ 2 \pi, & k = 0 \end{cases}
\end{align*}

Obținem, înlocuind în serie, că:
\[
    \ds\int_0^{2\pi} f(x) dx = a_0 \cdot 2\pi,
\]
de unde putem calcula primul coeficient imediat:
\[
    a_0 = \dfrac{1}{2 \pi} \int_0^{2\pi} f(x) dx.
\]

Din motive teoretice explicate la curs, de obicei, în calculul de mai sus se ia
coeficientul din fața integralei $ \dfrac{1}{\pi} $.

Prelucrări simple conduc apoi și la calculul celorlalți coeficienți
și avem $ b_0 = 0 $ și în rest:
\begin{align*}
    a_n &= \dfrac{1}{\pi} \int_0^{2\pi} f(x) \cos (nx) dx, \quad n \geq 1 \\
    b_n &= \dfrac{1}{\pi} \int_0^{2\pi} f(x) \sin (nx) dx, \quad n \geq 1.
\end{align*}

Mai general, dacă funcția are perioada principală $ T = 2 \ell $, atunci coeficienții se calculează cu:
\begin{align*}
    a_0 &= \dfrac{1}{\ell} \int_{-\ell}^\ell f(x) dx \\
    a_n &= \dfrac{1}{\ell} \int_{-\ell}^\ell f(x) \cos (nx) dx, \quad n \geq 1 \\
    b_n &= \dfrac{1}{\ell} \int_{-\ell}^\ell f(x) \sin (nx) dx, \quad n \geq 1.
\end{align*}

De asemenea, scurtături de calcul ne mai ajută să concluzionăm că:
\begin{itemize}
    \item Dacă $ f $ este impară, atunci toți coeficienții $ a_n = 0 $;
    \item Dacă $ f $ este pară, atunci toți coeficienții $ b_n = 0 $.
\end{itemize}

Amintim, de asemenea, din liceu, că:
\begin{itemize}
    \item Dacă $ f $ este o funcție pară, atunci $ \ds\int_{-a}^a f(x) dx = 2 \cdot \ds\int_0^a f(x) dx $;
    \item Dacă $ f $ este o funcție impară, atunci $ \ds\int_{-a}^a f(x) dx = 0 $.
\end{itemize}

Problema convergenței seriei Fourier este rezolvată de rezultatele de mai jos.

\begin{theorem}[G.\ L.\ Dirichlet]
    \label{thm:dirichlet}
    Fie $ f : \RR \to \RR $ o funcție periodică (de perioadă $ T = 2\pi $),
    mărginită, cu un număr finit de discontinuități de speța întîi și cu derivate
    laterale în orice punct.

    Atunci seria Fourier asociată lui $ f $ converge punctual în orice $ x \in \RR $ la
    \[
        \dfrac{1}{2} \left[ f(x + 0) + f(x - 0) \right].
    \]
    În particular, dacă $ f $ este continuă, atunci suma seriei Fourier este chiar
    funcția $ f(x) $ și avem:
    \[
        f(x) = \ds\sum_{n \geq 0} \left[ a_n \cos nx + b_n \sin nx \right], \quad \forall x \in \RR.
    \]
    
    Convergența uniformă se obține dacă $ f $ este de clasă $ \kal{C}^1 $ (măcar pe porțiuni)
    pe domeniul de definiție.
\end{theorem}

În exerciții, mai este de folos și \textbf{identitatea lui Parseval}:
\[
    \dfrac{|a_0|^2}{2} + \ds\sum_{n \geq 1} |a_n|^2 + |b_n|^2 = \dfrac{1}{\pi} \int_0^{2\pi} |f(x)|^2 dx.
\]

\vspace{0.3cm}

\textbf{Exemplu:} Fie funcția $ f(x) = |x| $, cu $ x \in [-\pi, \pi] $.
\begin{enumerate}[(a)]
    \item Calculați coeficienții Fourier și seria Fourier a lui $ f $, prelungind-o la $ \RR $.
    \item Calculați suma seriei $ \dfrac{1}{1^2} + \dfrac{1}{3^2} + \dfrac{1}{5^2} + \dots $.
\end{enumerate}

\emph{Soluție:} Prelungim funcția la $ \RR $ prin \qq{repetarea} acesteia pe toată axa reală.
Apoi, calculăm coeficienții, pe rînd:
\begin{itemize}
    \item $ a_0 = \dfrac{1}{\pi} \ds\int_{-\pi}^{\pi} |x| dx = \pi $;
    \item $ a_n = \dfrac{1}{\pi} \ds\int_{-\pi}^{\pi} |x| \cos nx dx = \dfrac{2}{\pi n^2} \left( (-1)^n - 1 \right) $;
    \item $ b_n = 0, \forall n \geq 0 $, deoarece $ f $ este impară.
\end{itemize}

Așadar, seria Fourier (care este convergentă, deoarece $ f $ este de clasă $ \kal{C}^1 $ pe porțiuni) este:
\[
    f(x) = |x| = \dfrac{\pi}{2} + \sum_{n \geq 1} \dfrac{2}{\pi n^2} \left[ (-1)^n - 1 \right] \cos nx.
\]

Remarcăm, în plus, că pentru $ n $ par, seria este nulă.

În plus, pentru $ n = 0 $, se obține exact suma seriei cerute, anume:
\[
    \dfrac{1}{1^2} + \dfrac{1}{3^2} + \dfrac{1}{5^2} + \dots = \dfrac{\pi^2}{8}.
\]



%%%%%%%%%%%%%%%%%%%%%%%%%%%%%%%%%%%%%%%%%%%%%%%%%%%%%%%%%%%%%%%%%%%%%%%%%%%%%%%%%%%%%%%%%%%%%%%%%%%%%%%%%%%%%%%%%%

\section{Exerciții}

Găsiți dezvoltarea în serie Fourier pentru funcțiile: % OLTEANU & COSTACHE
\begin{enumerate}[(1)]
    \item $ f(x) = x, x \in (-\pi, \pi) $.
    \item $ f(x) = x^2, x \in (-\pi, \pi] $. Calculați apoi suma $ \ds\sum_{n \geq 1} \dfrac{1}{n^2} $.
    \item $ f(x) = e^{ax}, x \in (-\pi, \pi] $, cu $ a \in \RR $. Calculați apoi suma $ \ds\sum_{n \geq 1} \dfrac{1}{n^2 + a^2} $.
    \item $ f(x) = \dfrac{\pi - x}{2}, x \in [0, 2\pi) $.
    \item $ f(x) = \pi^2 - x^2, x \in (-\pi, \pi ) $. Calculați apoi suma $ \ds\sum_{n \geq 1} \dfrac{1}{n^2} $.
    \item $ f(x) = | \sin x |, x \in (-\pi, \pi] $.
    \item $ f(x) = | \cos x |, x \in (-\pi, \pi) $.
\end{enumerate}